\section{Wire Format}

Two new ZAP schema IDs allocated for this LP, in the range
\texttt{0xF0}--\texttt{0xFF}. The range is chosen to avoid collision
with LP-201's \texttt{0xD0}--\texttt{0xDF} reservation and with any
future rollup-substrate allocation in the \texttt{0xE0}--\texttt{0xEF}
range that LP-204 leaves open (note: LP-218 takes
\texttt{0xE0}--\texttt{0xEF} for rollup envelopes; LP-208 takes
\texttt{0xF0}--\texttt{0xFF}; no overlap).

\subsection{Schema \texttt{0xF0}: DAGHeader}

Carried on unidirectional QUIC streams (LP-201) from emitter to
peers; replicated to the DHT (LP-201 layer~C) under
\texttt{sha256(header.Bytes())}.

\begin{table}[h]
\centering
\small
\begin{tabular}{@{}llp{6.5cm}@{}}
\toprule
\textbf{Field} & \textbf{Type} & \textbf{Notes} \\
\midrule
\texttt{version} & \texttt{uint8} & DAG protocol version; activation initializes to \texttt{1} \\
\texttt{validator} & \texttt{[20]byte} & NodeID of the emitter \\
\texttt{epoch} & \texttt{uint64} & DAG epoch; advances on validator-set change \\
\texttt{seq} & \texttt{uint64} & Per-validator monotonic sequence number; gaps allowed \\
\texttt{timestamp} & \texttt{int64} & Emitter wall-clock at emission (Unix nanoseconds) \\
\texttt{parent\_hashes} & \texttt{[]bytes32} & Parent header hashes, one per other validator; zero-padded if no parent available; length-prefixed list \\
\texttt{tx\_count} & \texttt{uint16} & Number of tx hashes in this header \\
\texttt{tx\_hashes} & \texttt{[tx\_count][32]byte} & List of tx hashes for txs admitted since previous header \\
\texttt{body\_hash} & \texttt{[32]byte} & \texttt{sha256(tx\_hashes\_bytes)}; separate from header hash to allow body GC without invalidating header references \\
\texttt{sig} & \texttt{[96]byte} & BLS sig over the canonical pre-image; matches LP-022 BLS leg curve \\
\bottomrule
\end{tabular}
\caption{DAGHeader (schema \texttt{0xF0}) wire layout.}
\label{tab:dagheader}
\end{table}

The BLS pre-image for the signature is
\texttt{version $\|$ validator $\|$ epoch $\|$ seq $\|$ timestamp $\|$
parent\_hashes $\|$ body\_hash}. The \texttt{body\_hash} is the
content-address of the DAGBody keyed in the DHT.

\subsection{Schema \texttt{0xF1}: DAGBody}

Carried on bidirectional QUIC streams (LP-201) on
\texttt{DHTFindValue} response; stored in the DHT under
\texttt{body\_hash}.

\begin{table}[h]
\centering
\small
\begin{tabular}{@{}llp{7cm}@{}}
\toprule
\textbf{Field} & \textbf{Type} & \textbf{Notes} \\
\midrule
\texttt{version} & \texttt{uint8} & DAG protocol version; matches the parent DAGHeader \\
\texttt{tx\_count} & \texttt{uint16} & Number of txs \\
\texttt{tx\_bytes} & \texttt{[tx\_count][]byte} & Length-prefixed list of full tx bytes -- each tx is itself a ZAP frame \\
\bottomrule
\end{tabular}
\caption{DAGBody (schema \texttt{0xF1}) wire layout.}
\label{tab:dagbody}
\end{table}

A DAGBody is content-equivalent to a list of tx bytes. There is no
header re-embedded in the body -- the body is a pure container so
that the same body bytes can serve multiple equivalent DAGHeaders
that happen to commit to the same tx-set.

\subsection{Schema ID allocation}

\begin{table}[h]
\centering
\small
\begin{tabular}{@{}lll@{}}
\toprule
\textbf{Schema ID} & \textbf{Type} & \textbf{Notes} \\
\midrule
\texttt{0xC0}--\texttt{0xCB} & chains-VM wire (12 VMs) & LP-186 \\
\texttt{0xD0}--\texttt{0xDF} & ZAP P2P transport & LP-201 \\
\texttt{0xE0}--\texttt{0xEF} & rollup-VM envelopes & LP-218 \\
\texttt{0xF0} & DAGHeader & This LP \\
\texttt{0xF1} & DAGBody & This LP \\
\texttt{0xF2}--\texttt{0xFF} & Reserved & DAG extensions (Mysticeti ordering hints in LP-209, Block-STM speculation hints in LP-210, cross-shard parent pointers in LP-211) \\
\bottomrule
\end{tabular}
\caption{Cross-LP schema ID allocation map.}
\label{tab:schema-ids}
\end{table}

Schema ID rationale: the original LP-208 draft suggested
\texttt{0xE0}--\texttt{0xE1} but those collide with the rollup
substrate reservation that LP-218 finalizes. This LP uses
\texttt{0xF0}--\texttt{0xFF} instead. The change is mandatory --
schema IDs cannot overlap across LPs.
